\documentclass[11pt]{article}
\usepackage[letterpaper,margin=0.8in]{geometry}
\usepackage{array,booktabs,xcolor}
\setlength{\parindent}{0pt}\setlength{\parskip}{6pt}
\renewcommand{\arraystretch}{1.5}
\begin{document}
{\LARGE\bfseries Unit 4, Day 10: review day run sheet}\par
Study guide for 4.0 to 4.7, which mirrors the test part by part. Files: \texttt{day-10-review-study-guide.pdf} and its key.\par\medskip
\begin{tabular}{@{}p{1.1in}p{2.4in}p{3.1in}@{}}
\textbf{Clock} & \textbf{Do} & \textbf{Watch for}\\\hline
0:00--5:00 & Hand out the study guide. Post the key at the back, face down. Students start part A. & Students who stop at the first blank. Tell them to write the frame and keep going.\\
5:00--35:00 & Students work through parts A to E at their own pace, then check each part against the key before moving on. Walk the room; pull a group together for any part many are missing. & Part C formulas (swapped or unreduced subscripts) and part D prefixes (mono- on the first element). Dots left in a metal ion's brackets.\\
35:00--45:00 & Whiteboard rounds on the two parts the room found hardest: one question each from the guide, hold up, then compare with the key. & A part everyone found easy: skip it.\\
45:00--50:00 & Students write on the back: the one part they will study tonight, and one thing they still get wrong. Collect, read after class. & Do not collect the guide itself; it is theirs.\\
\end{tabular}\par\medskip
\textbf{Where to reteach.} Use the day 9 slips and Quiz 1 for who needs what: trends (Day 1 board), dots and the octet (Day 2 board), ion charges, formulas and names (Days 4 and 5 boards), Lewis structures (Day 7 board), prefixes and bond types (Day 8 board).\par
\textbf{The test.} Tomorrow, 61 points, 19 questions; question 19 is the stretch item. A periodic table is provided.\par
\textbf{If you run short.} Drop the whiteboard rounds. Never drop the check against the key; it is how they find their gaps.
\end{document}
